\section{Introduction}
Cluster Editing asks for the fewest edge insertions and deletions that turn a graph $G$ into a disjoint union of cliques.
Every induced $P_3$ must be hit by an edit, so a packing $H$ of induced $P_3$'s that pairwise share no vertex pair (a \emph{modification-disjoint} packing) gives the lower bound $|H|$ on the solution size.
Li, Pilipczuk and Sorge~\cite{LPS23} parameterize by the \emph{excess} $\ell$ over this bound (problem CEaMP): is there a solution of size at most $|H|+\ell$?
They show NP-hardness already for $\ell=0$, with every vertex in at most $23$ packed $P_3$'s, and an $n^{2\ell+O(1)}$ algorithm when every vertex lies in at most two packed $P_3$'s; whether the latter case is fixed-parameter tractable is open~\cite{LPS23}.
For vertex-disjoint packings the problem is FPT in $\ell$~\cite{BFK18}.

The question of this paper is which measurable \emph{interaction structure} of $H$ makes the local repair choices (one edit per packed $P_3$, in three ways) combine into a global solution in $f(\ell,w)\,n^{O(1)}$ time.
We take $w=\tw(I(H))$, where the intersection graph $I(H)$ has the packed $P_3$'s as nodes and joins two of them when they share a vertex.

\paragraph{What is proved.}
For $\ell=0$ the problem is a constraint satisfaction problem (CSP) with one variable of domain size $3$ per packed $P_3$ and constraint scopes that are cliques of $I(H)$ (Section~\ref{sec:l0}); dynamic programming over a tree decomposition gives $3^{w+1}n^{O(1)}$.
For general $\ell$, guessing the at most $\ell$ edits outside all packed $P_3$'s reduces to the same CSP with domain size $7$, giving $7^{w+1}n^{2\ell+O(1)}$ (Section~\ref{sec:excess}).
Section~\ref{sec:structure} proves a strict twin lemma and, for packings $H$ that are maximal among pair-disjoint packings, the structure of the vertices $R$ outside $V(H)$ and the statement (C1) that in every optimum no edit has both endpoints in $R$; it derives frozen classes, a shrink rule whose threshold is tight, the bound $|R|\le(\ell+1)|V(H)|$, and a classification of the partners of edits touching $R$. A five-vertex example shows that the hypothesis cannot be weakened to maximality among vertex-disjoint packings, and a four-cycle shows that edits between $R$ and $V(H)$ cannot be excluded.
Section~\ref{sec:K} proves Theorem~K: with $s$ the maximum number of vertices spanned by a component of $I(H)$, CEaMP runs in $1.62^{(\binom s2+1)\ell}n^{O(1)}+n3^s$ time and has a kernel with at most $(s^2-s+2)\ell$ vertices; $I(H)$ a perfect matching is therefore FPT in $\ell$. The proof applies the reduction rule of~\cite{BFK18} to the components of $I(H)$ after proving its hypotheses in full; it still relies on their published proof of the rule's soundness, on a Cluster Editing solver and on the $2k$-vertex kernel as quoted there.
Section~\ref{sec:packing} proves counting lemmas for solutions with respect to an arbitrary packing (maximalisation with its cost in treewidth, a matching bound for hubs of deletions, small merges of pieces, and treewidth of the soft graph) and ends with a clearly labelled heuristic discussion of why simple hardness gadgets fail.
Section~\ref{sec:locality} proves that in an optimal solution every cluster has diameter at most $2$ and classifies every inserted unpacked pair into two types (T1: two soft steps through a common neighbour; T2: one kept hard step and one soft step). Pieces that share a cluster and have no soft pair between them have all cross pairs of type T1, which are local in a decomposition of $I(H)$ with at most $4c(c-1)$ partners per vertex; examples with $n=6,7,8$ show that merges of pieces joined by a soft pair can be of type T2, so the locality is proved only for soft-free merges. Section~\ref{sec:dproute} proves a normal form of optimal solutions for pair-maximal packings, a bound on the candidates of one endpoint of an edited hard pair given the other, an $n^{\ell}f(\ell,w)7^{w+\ell+1}$ algorithm, and an FPT algorithm in $(\ell,D,\omega)$ for the treewidth $\omega$ of the graph of edges and packed pairs of the non-\emph{big} part; a family with $\ell=w=0$ shows that $\omega$ is unbounded in $(\ell,w)$. These sections come with a status table and mark the missing reduction lemma as a conjecture.
Under ETH the exponent $\Theta(w)$ cannot be improved to $o(w)$ even at $\ell=0$; $w$ alone is insufficient; $\tw(G)$ and $\tw(I(H))$ are incomparable (Section~\ref{sec:limits}).

\paragraph{What is not proved.}
\textbf{Whether CEaMP is FPT in $\ell+\tw(I(H))$ (or W[1]-hard) is open}, already for $w=1$ with tree components of unbounded size. The algorithm of Section~\ref{sec:excess} is XP in $\ell$, Theorem~K needs bounded components, and no hardness is known. The normal form of Section~\ref{sec:structure} applies to a pair-maximal extension $H'$ of the given packing, whose intersection graph may have larger treewidth, so it does not by itself give a preprocessing for the parameter $(\ell,\tw(I(H)))$. What is left in precise form is: the pointers of type T2, the pointers realised by deleted hard edges, and the reduction lemma ML (Conjecture~\ref{conj:ML}), which would make Theorem~\ref{thm:omega} an FPT algorithm in $\ell+w$; whether merge pointers can be forced at unbounded distance in $I(H)$ is also open (Section~\ref{sec:left}). The algorithms of Sections~\ref{sec:excess} and~\ref{sec:dproute} remain XP in $\ell$. Section~\ref{sec:open} states the open problems.
